\begin{abstract}
This dissertation evaluates intrusion detection for RPL and 6LoWPAN networks under controlled changes in attack behaviour. A Contiki-NG and Cooja campaign implemented nine attack families: blackhole, sinkhole, DIS flooding, grayhole, increase-rank, DIO suppression, worst-parent, wormhole, and Sybil. Five attack and five matched control seeds were validated for each family, producing 90 baseline runs and a frozen dataset of 810 labelled 60-second routing windows. Static Gaussian and CART classifiers were evaluated across attack families using complete simulation seeds as the unit of separation. The CART model produced zero attack recall in 48 of 72 cross-attack pairs and mean F1 of 0.1822, demonstrating poor transfer when the attack mechanism changed. Whole-seed adaptation increased mean CART F1 to 0.8258 with one target seed and 0.8792 with three target seeds. Sybil identity manipulation was introduced as a new attack surface; static Sybil-target F1 was 0.2183, while one Sybil adaptation seed increased held-out F1 to 0.9683. A rank-state CUSUM monitor and delayed-error DDM both detected five held-out sinkhole shifts with no alert in five matched controls. Further experiments showed sensitivity to attacker relocation and lossy radio conditions. A direct rank-parent sinkhole intervention caused severe churn, reducing application responses by 79.06 percent and increasing radio events by 571.22 percent. The findings show that adaptation can recover performance when the feature representation exposes the changed mechanism, but retraining and mitigation are not interchangeable solutions. The contribution is a reproducible Cooja dataset and evaluation framework for cross-attack transfer, seed-safe adaptation, online drift detection, robustness, and critical defence analysis.
\end{abstract}
